%_________________
\section{Distribusi kontinu}
\label{contDist}

\noindent%
Sejauh ini dalam bab ini kita telah membahas kasus
ketika hasil suatu variabel bersifat diskret.
Dalam bagian ini, kita mempertimbangkan konteks ketika
hasilnya merupakan variabel numerik kontinu.

\index{data!US adult heights|(}
\index{hollow histogram|(}
\begin{examplewrap}
\begin{nexample}{Gambar~\ref{fdicHistograms} menampilkan beberapa
    histogram berongga yang berbeda untuk tinggi badan orang dewasa di AS.
    Bagaimana perubahan jumlah bin memungkinkan Anda membuat
    penafsiran yang berbeda terhadap data tersebut?}
  \label{usHeights}%
  Penambahan jumlah bin memberikan perincian yang lebih besar.
  Sampel ini sangat besar, sehingga bin yang jauh lebih sempit
  masih dapat digunakan dengan baik.
  Biasanya kita tidak menggunakan bin sebanyak ini pada sampel
  yang lebih kecil karena sedikitnya hitungan per bin membuat tinggi bin
  sangat tidak stabil.
\end{nexample}
\end{examplewrap}

\begin{figure}[ht]
  \centering
  \Figure[Ditampilkan empat histogram berongga tinggi badan orang dewasa di AS dalam sentimeter dengan lebar bin yang berbeda-beda. Perbedaan tampilannya dibahas terlebih dahulu, lalu bentuk histogram terakhir yang paling terperinci dijelaskan. Histogram pertama memiliki sekitar 6 bin dengan nilai yang tampak tidak nol, sehingga garis luarnya sangat menyerupai kotak-kotak. Histogram kedua memiliki sekitar 12 bin tidak nol dan tampak sedikit lebih kontinu serta tidak sekaku histogram pertama. Histogram ketiga memiliki sekitar 25 bin tidak nol dan garis luar histogram berongganya mulai tampak agak lebih mulus. Histogram terakhir memiliki sekitar 50 bin; karena jumlah bin yang besar, distribusinya tampak cukup mulus, yakni tidak ada langkah dari satu bin ke bin berikutnya yang berupa kenaikan atau penurunan tinggi yang besar. Berikutnya, histogram terakhir ini dijelaskan: tinggi bin kira-kira nol sampai 147, lalu naik secara bertahap hingga sekitar 155 sebelum meningkat tajam sedikit sampai 157, kemudian naik secara bertahap ke puncak sekitar 165. Dari sini histogram turun sekitar 10\% dari puncaknya pada 170, lalu penurunannya lebih bertahap sampai sekitar 183. Pada titik itu histogram menurun cepat sampai sekitar 187, kemudian mulai menurun lebih perlahan sambil mendekati 0. Pada sekitar 200, tinggi bin pada dasarnya telah mencapai nol dan tetap di sana.]{}{fdicHistograms}
  \caption{Empat histogram berongga tinggi badan orang dewasa di AS
      dengan lebar bin yang berbeda-beda.}
  \label{fdicHistograms}
\end{figure}

\begin{examplewrap}
\begin{nexample}{Berapa proporsi sampel yang memiliki tinggi antara \resp{180} cm dan \resp{185} cm (sekitar 5 kaki 11 inci sampai 6 kaki 1 inci)?}\label{contDistProb}
Kita dapat menjumlahkan tinggi bin-bin dalam rentang \resp{180} cm sampai \resp{185} cm lalu membaginya dengan ukuran sampel. Sebagai contoh, hal ini dapat dilakukan menggunakan dua bin berarsir pada Gambar~\ref{usHeightsHist180185}. Kedua bin dalam rentang ini memiliki hitungan 195,307 dan 156,239 orang, sehingga menghasilkan estimasi peluang berikut:
\begin{align*}
\frac{195307 + 156239}{\text{3,000,000}} = 0.1172
\end{align*}
Pecahan ini sama dengan proporsi luas histogram yang berada dalam rentang \resp{180} sampai \resp{185} cm.
\end{nexample}
\end{examplewrap}

\begin{figure}[h]
  \centering
  \Figure[Ditampilkan sebuah histogram tinggi badan dengan dua bin histogram antara 180 dan 185 sentimeter diberi arsiran; bin-bin tersebut merepresentasikan individu yang tingginya antara 180 dan 185 sentimeter.]{0.9}{usHeightsHist180185}
  \caption{Histogram dengan lebar bin 2.5 cm.
      Daerah berarsir merepresentasikan individu dengan
      tinggi antara \resp{180} dan \resp{185} cm.}
  \label{usHeightsHist180185}
\end{figure}


\D{\newpage}

\subsection{Dari histogram ke distribusi kontinu}

Perhatikan peralihan dari histogram berongga yang menyerupai kotak-kotak di bagian kiri atas Gambar~\ref{fdicHistograms} ke grafik yang jauh lebih mulus di bagian kanan bawah. Pada grafik terakhir ini, bin-bin begitu sempit sehingga histogram berongganya mulai menyerupai kurva mulus. Hal ini menunjukkan bahwa tinggi badan populasi sebagai variabel numerik \emph{kontinu} mungkin paling baik dijelaskan dengan kurva yang merepresentasikan garis luar bin-bin yang sangat sempit.

Kurva mulus ini merepresentasikan
\termsub{fungsi kepadatan peluang}
    {probability!density function}
(juga disebut \term{kepadatan} atau \term{distribusi}), dan kurva semacam itu ditampilkan pada Gambar~\ref{fdicHeightContDist} serta ditumpangkan pada histogram sampel. Suatu kepadatan memiliki sifat khusus: luas total di bawah kurva kepadatan adalah 1.

\begin{figure}[tbh]
\centering
\Figure[Ditampilkan sebuah histogram tinggi badan orang dewasa di AS dengan garis kontinu yang ditumpangkan mengikuti tinggi bin-bin. Garis kontinu ini mulus dan merepresentasikan bentuk histogram berongga seandainya kita memiliki data tak terhingga dan lebar bin begitu kecil sehingga garis luar histogram yang menyerupai kotak-kotak tampak kontinu. Bentuk ini disebut “kepadatan peluang kontinu”.]{0.9}{fdicHeightContDist}
\caption{Distribusi peluang kontinu tinggi badan orang dewasa di AS.}
\label{fdicHeightContDist}
\end{figure}

\index{hollow histogram|)}


\D{\newpage}

\subsection{Peluang dari distribusi kontinu}

Dalam Contoh~\ref{contDistProb}, kita menghitung proporsi individu dengan tinggi \resp{180} sampai \resp{185} cm sebagai pecahan:
\begin{align*}
\frac{\text{jumlah orang dengan tinggi antara \resp{180} dan \resp{185}}}{\text{ukuran sampel total}}
\end{align*}
Kita menemukan jumlah orang dengan tinggi antara \resp{180} dan \resp{185} cm dengan menentukan proporsi luas histogram dalam rentang tersebut. Serupa dengan itu, kita dapat menggunakan luas daerah berarsir di bawah kurva untuk menemukan peluang (dengan bantuan komputer):
\begin{align*}
P(\text{\var{height} antara \resp{180} dan \resp{185}})
	= \text{luas antara \resp{180} dan \resp{185}}
	= 0.1157
\end{align*}
Peluang bahwa orang yang dipilih secara acak memiliki tinggi antara \resp{180} dan \resp{185} cm adalah 0.1157. Nilai ini sangat dekat dengan estimasi dari Contoh~\ref{contDistProb}: 0.1172.

\begin{figure}[h]
  \centering
  \Figure[Ditampilkan sebuah kurva kepadatan tinggi badan dengan daerah antara 180 dan 185 sentimeter diberi arsiran.]{0.7}{fdicHeightContDistFilled}
  \caption{Kepadatan tinggi badan populasi orang dewasa di AS
      dengan luas antara 180 dan 185 cm diberi arsiran.
      Bandingkan grafik ini dengan Gambar~\ref{usHeightsHist180185}.}
  \label{fdicHeightContDistFilled}
\end{figure}

\begin{exercisewrap}
\begin{nexercise}
Tiga orang dewasa di AS dipilih secara acak. Peluang bahwa satu orang dewasa memiliki tinggi antara \resp{180} dan \resp{185} cm adalah 0.1157.\footnotemark\vspace{-1.5mm}
\begin{enumerate}
\setlength{\itemsep}{0mm}
\item[(a)] Berapa peluang bahwa ketiganya memiliki tinggi antara \resp{180} dan \resp{185} cm?
\item[(b)] Berapa peluang bahwa tidak satu pun memiliki tinggi antara \resp{180} dan \resp{185} cm?
\end{enumerate}
\end{nexercise}
\end{exercisewrap}
\footnotetext{Jawaban singkat:
  (a) $0.1157 \times 0.1157 \times 0.1157 = 0.0015$.
  (b) $(1-0.1157)^3 = 0.692$}

\begin{examplewrap}
\begin{nexample}{Berapa peluang bahwa orang yang dipilih secara acak memiliki tinggi \textbf{tepat} \resp{180}~cm? Anggap Anda dapat mengukur dengan sempurna.}
\label{probabilityOfExactly180cm}
Peluang ini adalah nol. Seseorang mungkin memiliki tinggi yang mendekati \resp{180} cm, tetapi tidak tepat \resp{180} cm. Hal ini juga masuk akal berdasarkan definisi peluang sebagai luas; tidak ada luas yang tercakup antara \resp{180}~cm dan \resp{180}~cm.
\end{nexample}
\end{examplewrap}

\begin{exercisewrap}
\begin{nexercise}
Misalkan tinggi seseorang dibulatkan ke sentimeter terdekat. Apakah ada peluang bahwa tinggi \textbf{terukur} seseorang yang dipilih secara acak adalah \resp{180} cm?\footnotemark
\end{nexercise}
\end{exercisewrap}
\footnotetext{Peluang ini bernilai positif. Siapa pun yang tingginya antara \resp{179.5} cm dan \resp{180.5} cm akan memiliki tinggi \emph{terukur} sebesar \resp{180} cm. Ini mungkin merupakan skenario yang lebih realistis dijumpai dalam praktik daripada Contoh~\ref{probabilityOfExactly180cm}.}

\index{data!US adult heights|)}


{\input{ch_probability/TeX/continuous_distributions.tex}}
